\printedpage{29}
% head: A. Grothendieck : Dérivateurs, Ch. XV, Édition des Universités de Montpellier II et Paris VII
% foot: Version 5 novembre 2010

I again assume the following conditions on $\cM$:
\msp{page 44}

\begin{enumerate}[nosep]
\item[\textit{a})] $\cM$ stable under small direct limits (and not only filtering ones).
\item[\textit{b})] $\cM$ endowed with a cardinal filtration $(\Filt^{\pi})_{\pi\geq\pi_0}$.
\item[\textit{c})] The objects of $\cM$ are accessible (\textbf{N.B.} It suffices to assume this for the objects of $\Filt^{\pi_0}$).
\end{enumerate}

Condition ($b$) implies that $\cM$ contains a small generating family by strict epimorphisms (e.g. a small subcategory $\cC$ of $\uFilt^{\pi_0} = \cM^{\pi_0}$ such that the inclusion $\cC \longrightarrow \uFilt^{\pi_0}$ is an equivalence of categories). By SGA 4 I 8.12.7, it follows that $\cM$ is stable under small inverse limits. There follows the

\textbf{Lemma 5.1} (under conditions ($a$), ($b$)). \textit{For every object $X$ of $\cM$, the set} $\Ss(X)$ \textit{of subobjects of $X$, endowed with the natural order relation ``of inclusion'', admits} inf \textit{and} sup \textit{of arbitrary size.}

inf, since these are expressed as inverse limits in $\cM$. Now the existence of inf implies that of sup, since $\sup_J = \inf_{J'}$, where $J'$ is the set of upper bounds of $J$ in $\Ss(X)$.

\textbf{Corollary 5.2.} \textit{For every arrow $f : Z \longrightarrow X$ in $\cM$, there exists, among the subobjects $X'$ of $X$ through which $f$ factors, a smallest one} % [sic] lesquelles -> lesquels, "sous-objets" is masculine
\msp{page 45}
\textit{(which we will denote} $\Img f$\textit{). The morphism $Z \longrightarrow \Img f$ is an epimorphism.}

This last assertion comes from this one:

\textbf{Lemma 5.3.} \textit{Let $\cM$ be a category in which kernels of double arrows exist. Let $f : Z \longrightarrow Y$ be an arrow of $\cM$. Then $g$ is an epimorphism if and only if $g$ does not factor through a subobject that is} strict, $Y_0$, \textit{of $Y$ different from $Y$.} % [sic] g -> f, the arrow has been named f

We are going to make an additional (practically innocuous) hypothesis on the $\Filt^{\pi}$:

\begin{enumerate}[nosep]
\item[\textit{d})] $\Filt^{\pi}$ is stable under direct limits indexed by index categories $I$ such that $\card \Fl(I) \leq \pi$\footnote[16]{This condition will be called ($d_0$) below.}. (It suffices to assume stability under initial object, simple amalgamated sums, and filtering $\varinjlim$ indexed by an ordered set $I$ of cardinality $\leq \pi$.) Moreover, if $f : Z \longrightarrow X$ is an arrow of $\cM$ such that $Z$ or $X$ is in $\Filt^{\pi}$, then $\Img f \in \Filt^{\pi}$.
\end{enumerate}

This last condition on $\Img f$ splits into two:
\msp{page 46}

\begin{enumerate}[nosep]
\item[$d_1$)] Every subobject of an object of $\Filt^{\pi}$ is in $\Filt^{\pi}$.
\item[$d_2$)] If $Z$ is in $\Filt^{\pi}$, then every quotient $Z'$ of $Z$ such that $\Img p = Z'$ (where $p : Z \longrightarrow Z'$ is the canonical morphism) is in $\Filt^{\pi}$.\footnote[17]{It would seem that we will not need ($d_1$), only ($d_2$). But in fact, ($d_1$) must be a consequence of ($d_2$), and ($d_2$) is a consequence of the first part of ($d$) (stability under finite $\varinjlim$) if every epimorphism in $\cM$ is strict, hence effective (if the $\Filt^{\pi}$ are stable under finite $\varprojlim$ as well \ldots).}
\end{enumerate}

\printedpage{30}
% head: A. Grothendieck : Dérivateurs, Ch. XV, Édition des Universités de Montpellier II et Paris VII
% foot: Version 5 novembre 2010

Of course, in practice it will always be true that $\Filt^{\pi}$ is stable under passage to subobjects and quotients of an object (not only the particular quotients at issue in ($d_2$). % [sic] closing parenthesis of "(pas seulement" missing, parentheses unbalanced as printed
Besides, when every bimorphism (\textit{i.e.} monomorphism \textit{and} epimorphism) of $\cM$ is an isomorphism, then ($d_2$) is indeed equivalent to stability of $\Filt^{\pi}$ under quotients of objects.

\textbf{Proposition 5.3.} \texttt{[Twice 5.3.]} \textit{Suppose conditions} (a), (b), (d) \textit{are satisfied for $\cM$. Let $\pi$ be a cardinal $\geq \pi_0$, and $X$ an object of $\cM$. Let} $\Ss_{\pi}(X)$ \textit{be the set of subobjects of $X$ that are in} $\Filt^{\pi}$\textit{. This set, ordered by the order relation induced by} $\Ss(X)$\textit{, is large relative to $\pi$ (and in particular filtering). If, moreover, filtering direct limits of monomorphisms in $\cM$ are monomorphisms}\footnote[18]{This is an \textit{exactness condition} on $\cM$, which I will denote as:\par
\noindent\circled{e}\quad Monomorphisms in $\cM$ are stable under filtering $\varinjlim$ in $\uFl(\cM)$.} \textit{(e.g. if filtering direct limits are left exact), or if the objects of} $\Filt^{\pi}$ \textit{are $\pi$-accessible, we have $\varinjlim\limits_{\alpha\in\Ss_{\pi}(X)} Z_{\alpha} \isoto X$}\footnote[19]{Say also that $\Ss_{\pi}(X)$ is $\pi$-adapted, and the direct system $(Z_{\alpha})_{\alpha}$ as well.}.

\msp{page 47}
\textsc{Proof}. Let $(Z_{\alpha})_{\alpha\in A}$ be a family of subobjects of $X$ such that $Z_{\alpha} \in \Filt^{\pi}$, and $\card A \leq \pi$. Let $Z = \coprod_{\alpha\in A} Z_{\alpha}$, we have $Z \in \Filt^{\pi}$ by ($c$). The smallest subobject of $X$ that is an upper bound for the $Z_{\alpha}$ is none other than $\Img (Z \longrightarrow X)$, which is in $\Filt^{\pi}$ by ($d_1$). Let us prove that $\varinjlim\limits_{\alpha\in\Ss_{\pi}(X)} Z_{\alpha} \isoto X$ is an isomorphism, when we know that it is a monomorphism. Now
\[
X \;=\; \varinjlim_{I} X_i \ ,
\]
where $(X_i)_{i\in I}$ is a filtering direct system, $I$ an ordered set large relative to $\pi$, the $X_i$ in $\Filt^{\pi}$. For every $i \in I$, let $X'_i$ be the image of $X_i$ in $X$. We therefore have $X'_i \in \Filt^{\pi}$ by ($d_2$). Using the fact that the $X_i \longrightarrow X'_i$ are epimorphisms, we see at once that $X = \varinjlim\limits_{I} X'_i$. Consider the map
\[
\begin{aligned}
I \quad &\longrightarrow \quad \Ss_{\pi}(X)\\
i \quad &\longmapsto \quad X'_i \ .
\end{aligned}
\]
This is an increasing map, let $J$ be its image, a filtering subset, large relative to $\pi$, of $\Ss_{\pi}(X)$, such that
\[
\varinjlim_{\alpha\in J} Z_{\alpha} \ \isoto \ X \ .
\]
\msp{page 48}
We therefore have
\[
\underbrace{\varinjlim_{\alpha\in J} Z_{\alpha}}_{\simeq X} \ \longrightarrow \ \varinjlim_{\alpha\in\Ss_{\pi}X} Z_{\alpha} \ \xrightarrow{\ \varphi\ } \ X \ ,
\]
the composite being the identity. If we know that $\varphi$ is a monomorphism, it follows that it is an isomorphism (a monomorphism having a section). Now $\varphi$ will be a monomorphism if every filtering $\varinjlim$ of monomorphisms is a monomorphism, and then we win.

\printedpage{31}
% head: A. Grothendieck : Dérivateurs, Ch. XV, Édition des Universités de Montpellier II et Paris VII
% foot: Version 5 novembre 2010

If we assume that the elements of $\Filt^{\pi}$ are $\pi$-accessible, we win as well, for we find:

\textbf{Corollary 5.4.} \textit{If the objects of} $\Filt^{\pi}$ \textit{are $\pi$-accessible (e.g. if $\pi \geq \pi_1$ and $\pi = \pi^{\pi_0}$, e.g. $\pi \geq \pi_1$ and $\pi = 2^{c}$, with $c \geq \pi_0$,} cf. (3.2) \textit{p. 16), then for every filtering subset $J$ of} $\Ss(X)$ \textit{large relative to $\pi$, such that $\varinjlim\limits_{\alpha\in J} Z_{\alpha} \isoto X$, $J$ is cofinal in} $\Ss_{\pi}(X)$ \textit{(hence $\varinjlim\limits_{\alpha\in J} Z_{\alpha} \isoto \varinjlim\limits_{\alpha\in\Ss(X)} Z_{\alpha}$).}

This is immediate.

\msp{page 49}
\textbf{Scholie 5.5.} When we make on $\cM$ the fairly innocuous exactness hypothesis that filtering $\varinjlim$ of monomorphisms are monomorphisms, \textit{or} restrict ourselves to cardinals $\pi$ such that the objects of $\Filt^{\pi}$ are $\pi$-accessible, we can \textit{canonize} the representation of an object $X$ as $\varinjlim\limits_{i\in I} X_i$, with $X_i \in \Filt^{\pi}$ and $I$ large relative to $\pi$, by taking $I = \Ss_{\pi}(X)$, and the corresponding direct system. We also have the following pleasant additional property: every filtering subset $J$ of $I$, of cardinality $\leq \pi$, admits a supremum $i_J$, and (if we assume that $\cM$ satisfies condition ($e$), page 46, note 18), we will have
\[
X_{i_J} \;=\; \varinjlim_{i\in J} X_i \ ;
\]
in other words, we have a representation of $X$ by a direct system that is (perfectly) \textit{$\pi$-adapted} (\textit{cf.} XIII, p. 85, 86).

\msp{page 50}
Now let
\[
M \;\subset\; \Ob\cM
\]
be a subset of $\Ob\cM$, stable under the isomorphism relation. If $M$ is $\pi$-accessible (according to condition $\mathrm{b}_{\pi}$ of XIII, p. 88), then putting
\[
M^{\pi} \;=\; M \cap \Filt^{\pi} \ ,
\]
$M^{\pi}$ determines $M$ by the following condition on $X \in \Ob\cM$.
\[
\begin{aligned}
&X \in M \iff \text{the subset } \Ss^{M}_{\pi}(X) \text{ of } \Ss_{\pi}(X) \text{ consisting of the } X_{\alpha} \in \Ss_{\pi}(X)\\
&\text{such that } X_{\alpha} \in M \ (\iff X_{\alpha} \in M^{\pi}) \text{ is cofinal in } \Ss_{\pi}(X).
\end{aligned}
\tag*{(5.1)}
\]
Note that when $X \in \Filt^{\pi}$, this criterion is trivially valid, since then $\Ss_{\pi}(X)$ admits $X$ as greatest element, and the cofinal subsets are those containing this final object. Thus (5.1) is interesting only for the $X \in \Ob\cM \backslash \Filt^{\pi}$. Note that for \textit{every} subset $M^{\pi}$ of $\Filt^{\pi}$, we can define a subset $M$ of $\Ob\cM$ by condition (5.1), and we will trivially have
\[
M^{\pi} \;=\; M \cap \Filt^{\pi} \ . \tag*{(5.2)}
\]
But probably, although $M$ satisfies (5.1), it is not for all that necessarily $\pi$-accessible. Indeed, for it to be so, it must be stable under $\varinjlim$
\printedpage{32}
% head: A. Grothendieck : Dérivateurs, Ch. XV, Édition des Universités de Montpellier II et Paris VII
% foot: Version 5 novembre 2010
\msp{page 51}
large relative to $\pi$ of elements of $M^{\pi}$, which implies that every direct factor of an element of $M^{\pi}$ is in $M$, hence also in $M^{\pi}$, according to the stability condition ($d$) on $\Filt^{\pi}$. This is a particular condition on the subset $M^{\pi}$ of $\Filt^{\pi}$. So, we will have to have a condition on $\cM$ stronger than the validity of (5.1), namely:
\[
\parbox{0.72\textwidth}{For every $X \in \Ob\cM$ and for every subset $J$ of $\Ss_{\pi}(X)$, such that
\begin{enumerate}[nosep]
\item[\textit{a})] $J$ is cofinal in $\Ss_{\pi}(X)$ and
\item[\textit{b})] If $J' \subset J$ is a filtering subset of $J$, of cardinality $\leq \pi$, then its supremum $\alpha_{J'}$ in $\Ss_{\pi}(X)$ belongs to $J$,
\end{enumerate}
we have $X \in M$ if and only if $J^{M} \overset{\text{def}}{=} J \cap \Ss^{M}_{\pi}(X)$ is cofinal in $J$ (or equivalently, in $\Ss_{\pi}(X)$)\,\footnotemark[20].}
\tag*{(5.3)}
\]
\footnotetext[20]{Instead of $J$ cofinal in $\Ss_{\pi}(X)$, it even suffices to have $\varinjlim\limits_{\alpha\in J} X_{\alpha} \isoto X$ -- the corresponding condition on $M$ thus becomes more demanding. But if the objects of $\Filt^{\pi}$ are $\pi$-accessible, this condition is in fact equivalent to the weaker condition.}

If we already admit the ``it suffices'' of (5.1), then it suffices in (5.3) to assume $X \in M$ and require that $J^{M}$ then be cofinal in $J$.

But note the
\msp{page 52}

\textbf{Proposition 5.6.}\footnote[21]{False as it stands, true if we assume that $M^{\pi}$ is stable under filtering $\varinjlim$.} \textit{Suppose $\pi$ chosen in such a way that the objects of} $\Filt^{\pi}$ \textit{are $\pi$-accessible (e.g. $\pi = 2^{c}$, with $c \geq \pi_1$,} cf. (3.2), \textit{p. 16). Then, for a subset $M$ of} $\Ob\cM$\textit{, $M$ is $\pi$-accessible if and only if it satisfies criterion} (5.3).

We have seen that this is necessary (tautologically, in view of proposition 5.3 and scholie 5.5, bottom of page 49). Let us prove that this is sufficient, by first establishing the

\textbf{Lemma 5.7.} \textit{Under the conditions of} 5.6 \textit{on $\pi$, and condition} (e) \textit{on $\cM$} (p. 46), \textit{for every $\pi$-adapted representation $X = \varinjlim\limits_{i\in I} X_i$ of an object $X$ of $\cM$, the subset $I'$ of those $i \in I$ such that $X_i \longrightarrow X$ is a monomorphism, is cofinal in $I$. Moreover, $I'$ is stable under suprema in $I$ of filtering subsets $J$ of $I'$ of cardinality $\leq \pi$.}

Let us show how this lemma 5.7 implies the ``it suffices'' of proposition 5.6. So let $X = \varinjlim\limits_{i\in I} X_i$ be a $\pi$-adapted representation of $X$ ($X_i \in \Filt^{\pi}$, $I$ large relative to $\pi$, \textit{etc.}). We must prove that, if 5.3 is valid, then $X \in M$ if and only if the set $I_M$ is cofinal in $I$.

\msp{page 53}
Assume $X \in M$. Consider the subset $I' \subset I$ described in the lemma. We have a map $I' \longrightarrow \Ss_{\pi}(X)$, namely $i \longmapsto X'_i = \Img (X_i \hookto X)$, and we have
\[
X \simeq \varinjlim_{i\in I} X_i \ \simeq \ \varinjlim_{i\in I'} X_i \ \simeq \ \varinjlim_{i\in I'} X'_i \ \simeq \ \varinjlim_{\alpha\in J} Z_{\alpha} \ ,
\]
where $J$ is the image of $I'$ in $\Ss_{\pi}(X)$. Since the objects of $\Filt^{\pi}$ are $\pi$-accessible, we conclude that $J$ is cofinal in $\Ss_{\pi}(X)$. By 5.3, it follows that the set $J_M$ of those $\alpha \in J$
\printedpage{33}
% head: A. Grothendieck : Dérivateurs, Ch. XV, Édition des Universités de Montpellier II et Paris VII
% foot: Version 5 novembre 2010
such that $Z_{\alpha} \in M$ is cofinal in $J$, hence $\neq \emptyset$. Thus $I'_M$, which is the inverse image of $J_M$ in $I'$, is in any case nonempty. Applying the same result while replacing $I'$ by the subset $I'_{\geq i'_0}$, we see that $I'_M$ is cofinal in $I'$, hence also in $I$, since $I'$ is cofinal in $I$; thus $I_M \supset I'_M$ is cofinal in $I$, \quad \textit{q.e.d.}

Conversely, suppose that $I_M$ is cofinal in $I$, and prove $X \in M$. The difficulty here comes from the fact that the representation $X = \varinjlim\limits_{i\in I_M} X_i$ is no longer adapted to $\pi$ -- $I_M$ is not necessarily stable in $I$ under suprema (in $I$) of filtering subsets in $I_M$, $I'$ of cardinality $\leq \pi$. In short, we must
\msp{page 54}
prove that $M$ is stable under $\varinjlim$ large relative to $\pi$ without more, and I now doubt that this follows from just the hypotheses made here, which only concern direct systems where the transition morphisms are monomorphisms.

Here, I think, is a counterexample. Let us take
\[
M \quad = \quad \parbox[t]{0.6\textwidth}{set of objects of $\cM$ isomorphic to a fixed object $X_0 \in \Filt^{\pi}$, such that $X_0$ has no endomorphism which is a monomorphism other than the identity.}
\tag*{(5.4)}
\]
Then condition 5.3 is satisfied. For:
\begin{enumerate}[nosep]
\item[\textit{a})] If $X \in M$, \textit{i.e.} $X \simeq X_0$, then $X \in \Filt^{\pi}$, hence $\Ss_{\pi}(X)$ admits $X$ as final object, and $J$, being cofinal, contains this element, hence $J_M$ (which also contains it) is cofinal.
\end{enumerate}
And:
\begin{enumerate}[nosep]
\item[\textit{b})] If $J_M$ is cofinal, hypothesis (5.4) implies that the set of transition morphisms between the $Z_{\alpha}$ ($\alpha \in J_M$) consists of isomorphisms, hence the direct limit $X$ of the $Z_{\alpha}$ in question is isomorphic to each of them, hence $X \simeq X_0$, \textit{i.e.} $X \in M$.
\end{enumerate}
Yet, $M$ is $\pi$-accessible only if it is stable under $\varinjlim$ large relative to $\pi$, which implies that it is stable under \textit{direct factors.} But hypothesis
\msp{page 55}
(5.4) in no way rules out that $X_0$ admits a projector $p$ ($p^{2} = p$) that is nontrivial. This will have an image $X'$ in $\cM$, which will be a direct factor of $X_0$ not isomorphic to $X_0$, \textit{i.e.} not in $M$.

These difficulties disappear if we assume that $M^{\pi}$ is stable under filtering $\varinjlim$ indexed by ordered sets $I$ of cardinality $\leq \pi$. Indeed, in this case, in the preceding proof, the subset $I_M$ is stable under suprema (in $I$) of subsets $J$ of $I_M$ that are filtering and of cardinality $\leq \pi$. Thus the direct system $(X_i)_{i\in I_M}$ is still $\pi$-adapted, and we can apply lemma 5.7 to it, to find a cofinal subset $I'_M \subset I_M$, such that the $X_i \longrightarrow X$ for $i \in I'_M$ are monomorphisms. The image $J$ of $I'_M$ in $\Ss_{\pi}(X)$ is then a subset of $\Ss_{\pi}(X)$ such that $\varinjlim\limits_{\alpha\in J} Z_{\alpha} = X$, hence (by the hypothesis made on $\pi$) cofinal in $\Ss_{\pi}(X)$ (5.4), with $Z_{\alpha} \in M^{\pi}$ for every $\alpha \in J$. Thus by 5.3 we have $X \in M$, \textit{q.e.d.}

\msp{page 56}
It remains to prove lemma 5.7. With the hypotheses and notation of the lemma, so let $i_0 \in I$, we must find $i' \in I'$ such that $i' \geq i_0$. Consider the image $X'_i$ of $X_i$ in $X$, and